% ═══════════════════════════════════════════════════════════════════════
%  Stratégie 2 --- XAUUSD niveaux x10. Chiffres : macros générées.
% ═══════════════════════════════════════════════════════════════════════
\section{Limites}
\label{sec:x10_limites}

\lettrine[lines=2, findent=2pt, nindent=0pt]{C}{ette} section énumère, sans
atténuation, ce que ce travail ne démontre pas. Sur un dossier négatif, elle
a une fonction particulière~: elle dit à quelles conditions le verdict
pourrait être contesté --- et, pour chacune, si la contestation tiendrait.

\subsection{Le spread de sélection est constant}

C'est la limite la plus lourde. Le modèle de coût de sélection applique
\XTenSpreadUsd{} à chaque aller-retour, quelle que soit l'heure et quelle que
soit l'année. Le flux du courtier a depuis permis de mesurer une médiane de
\XTenSpreadMeasuredMedian{} sur sa fenêtre disponible et de publier une
sensibilité horaire, sans re-sélection. La couverture antérieure à cette
fenêtre reste extrapolée~; le coût constant demeure donc une limite du modèle
de sélection, même si son effet est maintenant quantifié.

\emph{Cette limite pourrait-elle renverser le verdict~?} Non. La
section~\ref{sec:x10_couts} montre que la sensibilité au spread mesuré reste
négative, tout comme l'extrapolation indicative à coût nul. Le meilleur
modèle disponible déplace le chiffre, pas son signe.

\subsection{La provenance amont du parquet n'est pas documentée}

Le fichier minute de référence porte une convention de datation maintenant
testée et réconciliée, mais sa provenance amont n'est pas documentée dans le
dépôt~: fournisseur, chaîne d'export, corrections éventuelles et version du
jeu de données ne forment pas un manifeste auditable. La correction de
datation résout un défaut identifié~; elle ne remplace pas cette traçabilité.

Cette lacune ne fournit aucun argument pour déployer une stratégie perdante.
Elle limite en revanche la reproductibilité externe et interdit d'attribuer
tous les écarts de flux au seul moteur d'exécution.

\subsection{Le seuil d'accélération a un double rôle}

Le même curseur règle deux choses~: la sensibilité de l'armement --- combien
d'accélération il faut pour commencer à surveiller un niveau --- et le
basculement du \emph{reversal} --- combien de décélération il faut pour jouer
le retour. C'est l'axe le plus discriminant de la grille, et il n'est donc
\textbf{pas interprétable scénario par scénario}~: on ne peut pas dire si les
valeurs élevées aident les \emph{breakouts}, aident les \emph{reversals}, ou
n'aident que parce qu'elles font trader moins.

Ce couplage était connu et publié avant la campagne, assumé pour tenir un
budget de vingt-sept configurations. Le découpler reviendrait à ajouter un axe
à la grille, ce que le gel interdit.

\subsection{La distance du stop n'agit que sur la moitié des scénarios}

Le troisième axe de la grille règle le stop du \emph{breakout}. Le stop du
\emph{reversal}, lui, est une constante gelée --- l'extrême du débordement
augmenté d'un demi-\code{ATR}. Un tiers de la grille ne concerne donc que la
moitié des scénarios, et c'est exactement l'axe qui ne classe rien dans les
résultats. Ce n'est pas une découverte de la campagne~: c'était une limite
annoncée de la grille.

\subsection{L'amorçage de la moyenne horaire}

La moyenne mobile exponentielle à 50~périodes horaires exige 50~barres closes.
Elle est donc indéfinie sur les deux premiers jours de l'échantillon, où le
contexte vaut zéro et où \textbf{aucun \emph{breakout} n'est possible} --- les
\emph{reversals}, eux, ne sont pas empêchés. L'effet est borné à deux séances
sur \XTenSampleSessions~: négligeable, mais il fallait le dire, parce qu'il
constitue une asymétrie de traitement entre les deux familles de scénarios sur
cette amorce.

\subsection{La durée maximale coupe des positions gagnantes}

\XTenTimeShare{} des transactions sortent à la 48\textsuperscript{e} barre,
avec une espérance de \metric{\XTenTimeR{}~R}. La durée de vie maximale gelée
tranche donc, en moyenne, des positions qui gagnaient.

\emph{Cette limite pourrait-elle renverser le verdict~?} On ne le sait pas, et
c'est précisément le problème~: allonger la durée maximale au vu de ce chiffre
serait un paramètre choisi après mesure, donc une sélection interdite. C'est
un point de conception à rouvrir si le mandat l'est
(section~\ref{sec:x10_non_deploiement}), sous une nouvelle spécification.

\subsection{Le Sharpe porte sur les rendements d'une stratégie souvent plate}

Les ratios publiés sont calculés sur des rendements \emph{de séance}. Or, sur
\XTenSampleSessions{} séances, une part importante ne porte aucune
transaction et vaut structurellement zéro~; et parmi celles qui en portent,
les séances à transaction unique dominent.

Un Sharpe annualisé sur cette base est une \textbf{convention de comparaison},
pas une propriété de la stratégie. Il permet de classer les vingt-sept
configurations entre elles --- ce pour quoi il est employé ici --- et il ne se
compare pas au Sharpe d'une stratégie continuellement exposée. Les séances
plates sont volontairement conservées dans l'index~: les retirer gonflerait le
ratio sans changer un dollar du résultat.

\subsection{Deux tests statistiques ne sont pas publiés}

Le test de supériorité prédictive du dépôt a une orientation inversée et rend,
sur une grille intégralement perdante, une valeur qui signifie l'inverse de ce
qu'elle paraît dire. Le test descendant de Romano--Wolf, lui, échoue
carrément. Les deux sont des limites préexistantes du module de tests
statistiques, signalées et \textbf{non corrigées} par cette campagne, et
aucun critère de la table de décision n'en dépend.

C'est une limite de \emph{couverture}~: le dossier aurait été plus complet
avec ces deux tests. Ce n'est pas une limite d'\emph{interprétation}~: leur
absence ne change rien au sens des
\XTenCriteriaFailingInterpretable{} critères en échec mesurés.

\subsection{Le portage d'exécution repose sur un modèle interpolé}

Le testeur MetaTrader~5 tourne en mode « OHLC M1 », qui interpole les
remplissages à l'intérieur de la minute. Les chiffres qu'il produit dépendent
donc d'un chemin intraminute reconstruit. Selon l'ordre des contacts, cette
interpolation peut améliorer ou dégrader le P\&L~: elle ne fournit aucune
garantie de borne supérieure ou inférieure par rapport à un compte réel.

S'ajoute une limite de couverture~: l'historique de l'or chez le courtier ne
remonte qu'à la fin 2022. Un peu plus de la moitié de l'échantillon --- et
précisément la moitié où le coût pèse le plus lourd --- ne sera jamais
vérifiable sur le moteur d'exécution.

\subsection{La dernière minute de séance diffère dans le portage QuantConnect}

La spécification et la référence considèrent comme dernière barre celle qui
débute à 16:57, heure de New York. Le portage QuantConnect conserve
\code{SESSION\_LAST\_MINUTE = 16:58}. Cette borne fonctionne dans le cloud,
avec la convention temporelle de LEAN, mais ne se reproduit pas telle quelle
sur l'export local redaté à l'ouverture.

L'écart est documenté et ne justifie aucune nouvelle optimisation. Il empêche
cependant de qualifier le portage complet d'identique à la référence tant que
la convention n'est pas rendue explicite des deux côtés.

\subsection{Une fenêtre horaire interdite, écart au mandat}

Aucune entrée n'est prise entre 16:30 et 18:15, heure de New York, autour de
la coupure quotidienne. Le mandat demandait la séance entière~: la stratégie
ne trade donc pas la réouverture asiatique. L'écart est inscrit dans la
spécification et son volume est faible --- c'est la plus petite des raisons
d'annulation de l'automate --- mais il existe.

\subsection{Aucun filtre d'annonces}

La stratégie ne connaît pas le calendrier économique. Une publication de
statistique américaine produit sur l'or des mouvements qui déclenchent
mécaniquement vitesse, accélération et franchissement de niveau. Le mandat
demandait la séance entière sans calendrier~; ce choix est respecté, et il
n'est pas mesuré séparément.

\subsection{Ce que ce rapport ne prétend pas établir}

\begin{insightbox}
\textbf{Ce document ne démontre pas qu'aucune stratégie de niveaux ronds ne
peut fonctionner sur l'or.} Il démontre que \emph{celle-ci} --- telle que le
mandat la définit, telle que la spécification la fige, dans les vingt-sept
réglages autorisés et sur sept années --- perd de l'argent, à toutes les
coupes examinées, et que cette perte survit au rééchantillonnage.

Une famille de signaux différente, une définition adaptative de la grille, ou
une restriction de séance sont des \emph{hypothèses non testées}. La
section~\ref{sec:x10_non_deploiement} les énonce comme telles.
\end{insightbox}

Enfin, ce rapport n'établit rien sur la \textbf{corrélation} de cette
stratégie au portefeuille existant du mandant. La question ne se pose pas~: un
actif dont l'espérance est négative n'améliore pas une allocation, quelle que
soit sa corrélation.
